% Generated draft: requires mathematical and source review.
\documentclass[11pt]{article}
\usepackage[margin=25mm]{geometry}
\usepackage{fontspec}
\setmainfont{Latin Modern Roman}
\usepackage{amsmath,amssymb,mathtools,unicode-math}
\setmathfont{Latin Modern Math}
\usepackage{enumitem,needspace,xurl,hyperref}
\setlength{\parindent}{0pt}
\setlength{\parskip}{5pt}
\setlength{\emergencystretch}{3em}
\newcommand{\N}{\mathbb{N}}
\newcommand{\Z}{\mathbb{Z}}
\newcommand{\Q}{\mathbb{Q}}
\newcommand{\R}{\mathbb{R}}
\newcommand{\C}{\mathbb{C}}
\newcommand{\F}{\mathbb{F}}
\newcommand{\A}{\mathbb{A}}
\newcommand{\PP}{\mathbb{P}}
\newcommand{\E}{\mathbb{E}}
\newcommand{\eps}{\varepsilon}
\newcommand{\dd}{\,\mathrm{d}}
\begin{document}
\section*{UnsolvedMath ID: 30005073}
\textbf{Source problem code:} catalogue entry
\subsection{Definitions and mathematical statement}
Let \(Y\) be a rational homology \(3\)-sphere, meaning that \(H_i(Y;\Q)\cong H_i(S^3;\Q)\) for every \(i\). The quantum invariants below are understood in their usual setting of closed, connected, oriented \(3\)-manifolds; the supplied excerpt does not specify their orientation and normalization conventions. In this setting the source's hypothesis \(b_1(Y)=0\), where \(b_1(Y)=\dim_{\Q}H_1(Y;\Q)\), is equivalent to the rational homology sphere condition. Write \(\operatorname{Spin}^c(Y)\) for the set of isomorphism classes of \(\operatorname{Spin}^c\) structures on \(Y\).

Define the Habiro ring by
\[
\widehat{\Z[q]}=\varprojlim_{n\geq1}\Z[q]\big/\left(\prod_{k=1}^{n}(1-q^k)\right).
\]
For a root of unity \(\zeta\), evaluation induces a well-defined homomorphism \(\operatorname{ev}_{\zeta}:\widehat{\Z[q]}\to\Z[\zeta]\): the defining product vanishes at \(\zeta\) once \(n\) is at least its order.

Put
\[
A_Y=H^1(Y;\C/2\Z),\qquad B_Y=H^1(Y;\Z/2\Z)\subseteq A_Y,
\]
using the inclusion of coefficient groups. The source uses the mod-\(2\) cohomological refinement \(\operatorname{WRT}_r(Y,\omega)\) for \(\omega\in B_Y\), and the Costantino--Geer--Patureau-Mirand invariant \(N_r(Y,\omega)\) for \(\omega\in A_Y\setminus B_Y\), at admissible levels \(r\). Let \(\mathcal T(Y,\omega)\) denote the version of Reidemeister--Turaev torsion required by the source. Retain its explicit hypothesis
\[
\mathcal T(Y,\omega)\ne0\qquad\text{for every }\omega\in A_Y\setminus B_Y.
\]
Subject to the admissibility and normalization conventions of the source, define
\[
\widetilde N_r(Y,\omega)=
\begin{cases}
\operatorname{WRT}_r(Y,\omega),&\omega\in B_Y,\\[2pt]
\displaystyle\frac{N_r(Y,\omega)}{\mathcal T(Y,\omega)},&\omega\in A_Y\setminus B_Y.
\end{cases}
\]
The particular torsion, the invariant normalizations, and the admissible levels are not defined in the supplied excerpt.

\textbf{Habiro-ring existence conjecture (source Conjecture 2).} For every such \(Y\), subject to the stated torsion hypothesis, there exists a family
\[
\bigl(I_s(Y)\bigr)_{s\in\operatorname{Spin}^c(Y)}\in
\prod_{s\in\operatorname{Spin}^c(Y)}\widehat{\Z[q]}
\]
such that, simultaneously for every source-admissible level \(r\) and cohomology class \(\omega\in A_Y\),
\[
\widetilde N_r(Y,\omega)=
\sum_{s\in\operatorname{Spin}^c(Y)}
C_s(Y,\omega)\,
\operatorname{ev}_{e^{2\pi i/r}}\bigl(I_s(Y)\bigr).
\]
The family depends on \(Y\) and \(s\), not on \(r\) or \(\omega\). The coefficients \(C_s(Y,\omega)\) are prescribed universal coefficients from the source, with dependence on \(r\) suppressed in its notation; they also involve certain homotopy-type invariants of \(Y\). They are not additional unknowns to be chosen to satisfy the equation. Their formula and the full admissibility conditions are absent, so this is a qualified reconstruction rather than a self-contained specification of those data.

The source motivates this conjecture through its preceding Conjecture 1, attributed there to references [3, 2], asking for series \(\widehat Z_s(Y)\in\Z[[q]]\), convergent for \(|q|<1\), whose weighted limits at \(q=e^{2\pi i/r}\) recover \(\widetilde N_r\). It states that root-of-unity limits alone do not ensure uniqueness of such series, because nonzero series may have zero limits at all roots of unity. The displayed Habiro-ring conjecture explicitly asserts existence; a separate uniqueness assertion for the indexed family is not supplied. The excerpt reports verification of the preceding conjecture for an unspecified family in reference [2], a reduction to an unspecified analogous conjecture for links in \(S^3\), and a possible analogous conjecture for \(\mathfrak{sl}(2|1)\) versions of CGP invariants in reference [5]. It also says that the torsion hypothesis holds for a generic \(3\)-manifold, asks how it might be avoided, notes that its conventions differ slightly from [2], and mentions a possible version for \(b_1(Y)>0\). None of these unspecified extensions or reported results supplies missing hypotheses for the statement above.
\subsection{Short English statement}
For each rational homology \(3\)-sphere satisfying the source's torsion nonvanishing hypothesis, seek one Habiro-ring element for each \(\operatorname{Spin}^c\) structure. Their evaluations at roots of unity should recover the combined refined WRT and torsion-normalized CGP invariant through the source's prescribed coefficients, for every admissible level and cohomology class. The coefficient formulas, admissibility rules, and normalization conventions remain unspecified. This is a statement-only draft requiring mathematical review.
\subsection{Sources}
\begin{itemize}
\item Local UnsolvedMath registry: \texttt{data/unsolvedmath/problems.json}, supplied record with permanent ID \(30005073\), source problem code \texttt{catalogue entry}, title ``Habiro-Ring Refinements of Three-Manifold Quantum Invariants.'' No individual proposer is supplied; the proposed year is recorded as 2022.
\item Supplied report citation: \emph{Mini-Workshop: Recent Developments in Representation Theory and Mathematical Physics (2023)}, Oberwolfach Reports, pp.~817--848, DOI \url{https://doi.org/10.4171/owr/2022/16}. The supplied literature list also labels this link ``Oberwolfach Report 2022/16.'' The title year and report identifier differ and have not been reconciled here.
\item Supplied literature reference: Costantino--Gukov--Putrov, SIGMA, \url{https://doi.org/10.3842/SIGMA.2023.010}. The supplied assessment attributes proofs of the \(q\)-series/root-of-unity relation under technical hypotheses and for several families to these authors.
\item Literature assessment checked at \texttt{2026-08-22}: this exact date records supplied dated curation, not independent verification. The registry labels the problem ``partially solved'' and states that the general existence and uniqueness conjecture for all rational homology spheres remains unproved. The supplied background explicitly requires independent verification of time-sensitive status and cited claims. The assessment's reference to uniqueness is broader than the explicit existence assertion in the excerpt.
\item The supplied excerpt attributes the preceding conjecture to [3, 2], CGP invariants to Costantino, Geer, and Patureau-Mirand [1], the mod-\(2\) WRT refinement to [6], and the \(\mathfrak{sl}(2|1)\) analogues to [5]. Full bibliographic entries and a reliable mapping of these numbered references to the supplied links are absent.
\item The registry records \texttt{statement\_status}: \texttt{corrected\_verified} and says the reconstruction was checked against the official report. These are supplied provenance claims; no source was visited or independently checked in preparing this draft.
\end{itemize}
\end{document}
